\documentclass[11pt]{article}
\usepackage[margin=2.5cm]{geometry}
\usepackage{amsmath,amsthm}
\usepackage[hidelinks]{hyperref}
\usepackage[paper=egg]{ledgercheck}
\newtheorem{theorem}{Theorem}
\newtheorem{lemma}[theorem]{Lemma}
\newtheorem{observation}[theorem]{Observation}
\newtheorem{conjecture}[theorem]{Conjecture}

\title{How long should I boil my egg?\\\large A worked example of an evidence ledger\ledgertitlenote}
\author{A.~Cook \and B.~Taster}
\date{}

\begin{document}
\maketitle

\begin{abstract}
A hard-boiled egg is done when its centre reaches $85\,^\circ$C. Starting from the cooking-time formula of
Buay et al., we show that bigger eggs need quadratically longer, that colder eggs need longer, and that an egg
taken from the fridge needs about one minute more than one at room temperature, however hard it should be.
A kitchen measurement agrees. This paper is an example for the trust-ledger tool: the kitchen, its data and
its authors are fictitious, and the appendix shows what the ledger records about each result.
\end{abstract}

\section{The model}

An egg is a prolate spheroid with semi-axes $a > b$. For heat flow, Buay et al.\ replace it by the sphere with
the same ratio of surface to volume, of radius
\begin{equation}\label{eq:radius}
  R = \frac{2ab}{b + \beta a}, \qquad \beta = \frac{\arcsin e}{e}, \quad e = \sqrt{1 - b^2/a^2}
\end{equation}
\cite[eq.~(22)]{Buay2006}, and they keep only the first term of the series solution for the temperature at its
centre. We take their result as given.

\begin{lemma}[Buay et al.]\label{lem:firstorder}
  A sphere of radius $R$ and thermal diffusivity $\alpha$, at temperature $T_0$, put into a bath at $T_b$,
  reaches the temperature $T_c$ at its centre, $T_0 < T_c < T_b$, after the time
  \[ t(R, \alpha, T_b, T_0, T_c) = \frac{R^2}{\pi^2 \alpha} \ln \frac{2 (T_b - T_0)}{T_b - T_c}. \]
\end{lemma}

The formula is their eq.~(11). It is an approximation (the first term of a series), best when the egg is
nearly done ($T_c$ close to $T_b$) and worst for a very soft egg; they compare it with a measurement and
find agreement within a few percent at $T_c = 85\,^\circ$C \cite{Buay2006}. An older and simpler estimate
\cite{Roura2000} is about 15\,\% too long.

\section{Three consequences}

\begin{lemma}[Bigger eggs]\label{lem:scaling}
  For every $k$, $t(kR, \alpha, T_b, T_0, T_c) = k^2\, t(R, \alpha, T_b, T_0, T_c)$.
\end{lemma}

\begin{lemma}[Colder eggs]\label{lem:colder}
  If $R, \alpha > 0$ and $T_0' < T_0 < T_c < T_b$, then
  \[ t(R, \alpha, T_b, T_0, T_c) < t(R, \alpha, T_b, T_0', T_c): \]
  a colder egg needs strictly longer.
\end{lemma}

\begin{theorem}[The fridge penalty]\label{thm:fridge}
  Assume Lemma~\ref{lem:firstorder}. If $R, \alpha > 0$ and $T_0' < T_0 < T_c < T_b$, an egg started at $T_0'$
  needs longer than one started at $T_0$ by
  \[ \frac{R^2}{\pi^2 \alpha} \ln \frac{T_b - T_0'}{T_b - T_0}, \]
  which does not depend on $T_c$.
\end{theorem}

For a medium egg ($R = 23.2$\,mm) and $\alpha = 1.6 \times 10^{-7}\,\mathrm{m^2/s}$ the prefactor is
$R^2/(\pi^2\alpha) = 341$\,s. In boiling water, an egg from the fridge ($5\,^\circ$C) instead of from the
counter ($20\,^\circ$C) needs $341\,\mathrm{s} \cdot \ln(95/80) \approx 59$\,s more: one minute, for a soft
egg as for a hard one.

\begin{lemma}[On a mountain]\label{lem:altitude}
  Where water boils below $100\,^\circ$C, every egg takes longer.
\end{lemma}

This lemma follows from the same formula, but we did not put it into the ledger: the appendix lists it as
not covered, and nothing is claimed about it there.

\section{In the kitchen}

We measured one egg: its semi-axes with a calliper, then its centre temperature with a probe thermometer every
30\,s after putting it, straight from the fridge, into boiling water. A script fits $\alpha$ to the readings
with the full series solution \cite[eq.~(5)]{Buay2006}.

\begin{observation}\label{obs:kitchen}
  Our egg ($a = 28.0$\,mm, $b = 21.5$\,mm, so $R = 23.2$\,mm by \eqref{eq:radius}), started at $5\,^\circ$C,
  reached $85\,^\circ$C at its centre after $14.4$\,min. The fitted thermal diffusivity is
  $\alpha = 1.60 \times 10^{-7}\,\mathrm{m^2/s}$, inside the range $1.5$ to $1.8 \times 10^{-7}\,\mathrm{m^2/s}$
  found by Buay et al.
\end{observation}

Three people ate the egg. In a real paper their notes would stay private, and the ledger would hold only a
salted commitment to them, so that they cannot be changed later without anyone noticing. In this example
the notes and the salt are in the supplement, so that the reader can check the commitment.

\begin{observation}\label{obs:taste}
  All three tasters called the egg firm but not rubbery, and none saw a green rim around the yolk.
\end{observation}

\begin{conjecture}\label{conj:six}
  Six minutes from boiling water makes the perfect soft-boiled egg.
\end{conjecture}

\section{What the reader can check}

\ledgersummary[What the ledger says]

\bibliographystyle{plain}
\bibliography{refs}
\end{document}
